\chapter{Einleitung}
\section{KI und das Fahrrad im Alltag}
\label{ch:intro}
Das Fahrrad ist aus dem heutigen Alltag nicht mehr wegzudenken, vor allem in der Stadt. In Zeiten der steigenden Benzinpreise, wachsenden Stadtbevölkerung und steigenden Unterhaltskosten sind die Menschen mehr denn je auf das Fahrrad angewiesen. Umso verwerflicher ist der fehlende Support in vielen Karten- und Routenprogrammen, vor allem wenn es um die Berücksichtigung vom entscheidenden Faktor der Straßenqualität geht.

Plattformen wie OpenStreetMap bieten zwar unterschiedliche Modelle zur Routenberechnung an, doch nur eines von ihnen bezieht die Straßenqualität in ihre Wegfindung wirklich mit ein. Zwar gibt es das plattformeigene Qualitätsmerkmal "smoothness" um über Straßenqualität zu entscheiden, doch sind diese Einordnungen subjektiv und werden erst nach dem Hochladen reviewt. Es ist diese Subjektivität, weswegen nicht alle Routenberechnungsmodelle „smoothness“ verwenden. Nur das Graphhopper-Modell hat es seit dem 4.0 release standardmäßig für Fahrräder in Benutzung.\footnote{\cite{Andreas2021}}

Diese Methode ist natürlich sehr zeitintensiv. Die benötigten Daten müssen händisch aufgenommen und mit den zugehörigen Tags hochgeladen werden. Danach werden sie erst auf Unstimmigkeiten reviewt. \footnote{{$https://wiki.openstreetmap.org/wiki/Contribute_map_data$}} Diese Methode benötigt dazu natürlich auch Einarbeitungszeit auf der Seite der Benutzer. Doch mithilfe von heutigen Machine Learning Methoden sind wir nicht mehr exklusiv auf subjektive Einschätzungen angewiesen. Es besteht die Möglichkeit, auch große Mengen an Straßendaten aufzunehmen und auszuwerten - im besten Fall auch ohne den Belag selbst mühselig untersuchen zu müssen.

Machine Learning ist ein Anwendungsbereich des Felds der Künstlichen Intelligenz. Diese Algorithmen erarbeiten eigenständig Muster und Features aus zugeführten Daten um aus ihenen zu lernen und vergleichbare Muster in neuen Daten wiederzuerkennen. Für die längste Zeit waren diese Anwendungen wegen mangelnder Hardwareleistung für den durchschnittlichen Benutzer unbrauchbar, doch mit den Entwicklungen in diesem Feld über die letzten Jahre, und den voranschreitenden Kapazitäten handelsüblicher Hardware lässt sich ein solcher Algorithmus sogar auf einem durchschnittlichen Smartphone betreiben. Die Grundlage für ein crowdbasiertes Sammeln und Klassifizieren von Daten ist heutzutage problemlos gegeben.
\section{Zielsetzung}

Das Ziel dieser Arbeit ist es, eine Grundlage für ein einfach zu erlernendes und gruppentaugliches Sammeln und objektiveres Klassifizieren von Daten zur Straßenqualität für Fahrradfahrer zu schaffen, das keine längeren Review-Prozesse benötigt. Hierzu werden mehrere Classifier präsentiert, der mithilfe von Smartphone-Sensordaten unterschiedliche Arten von Bodenbelägen und deren Qualität ermittelt, sowie eine Übersicht über die gesammelten Daten ausgeben. Bei der Qualtitätseinschätzung wird sich an den internationalen Standards für Oberflächenqualtiät wie dem IRI orientiert. Das Modell soll dazu leicht um neue Daten und Features erweiterbar sein.

\gls{library} as an example.

\subsection{And an even more important subsection}